% =============================================================================
% APENDICE: PARAMETROS OPTIMOS DE UMBRAL
% =============================================================================

\section{Par\'ametros \'Optimos de Umbral}
\label{app:optimal_params}

Esta secci\'on presenta los par\'ametros \'optimos de posicionamiento por percentil para cada modelo, obtenidos mediante grid search sobre el conjunto de validaci\'on.

\subsection{Metodolog\'ia de Optimizaci\'on}
\label{app:opt_methodology}

\subsubsection{Funci\'on Objetivo}

Optimizamos el Sharpe Ratio neto (despu\'es de costos) sobre el per\'iodo de validaci\'on:

\begin{equation}
(q_{3x}^*, q_{1x}^*) = \arg\max_{q_{3x}, q_{1x}} \text{Sharpe}(R_{net}(q_{3x}, q_{1x}))
\end{equation}

sujeto a $q_{3x} < q_{1x}$ (percentil para 3x debe ser m\'as selectivo que para 1x).

\subsubsection{Grid de B\'usqueda}

\begin{itemize}
    \item $q_{3x} \in \{5, 10, 15, 20, 25, 30\}$ (percentil superior para posici\'on UPRO)
    \item $q_{1x} \in \{20, 30, 40, 50, 60\}$ (percentil superior para posici\'on SPY)
    \item Total: 30 combinaciones evaluadas por modelo
\end{itemize}

\subsubsection{Restricciones}

\begin{enumerate}
    \item $q_{3x} < q_{1x}$: Asegura que UPRO (3x) requiera mayor confianza que SPY (1x)
    \item Ventana rolling: 63 d\'ias para c\'alculo de percentiles
    \item Warmup: Predicciones de training para los primeros 63 d\'ias de test
\end{enumerate}

\subsection{Tabla de Par\'ametros \'Optimos}
\label{app:optimal_table}

\begin{table}[htbp]
\centering
\caption{Par\'ametros \'Optimos de Umbral por Modelo}
\label{tab:optimal_params}
\small
\begin{tabular}{@{}llrrr@{}}
\toprule
\textbf{Modelo} & \textbf{Categor\'ia} & \textbf{$q_{3x}$ (\%)} & \textbf{$q_{1x}$ (\%)} & \textbf{Interpretaci\'on} \\
\midrule
Ridge & ML & 20 & 30 & Top 20\% $\rightarrow$ 3x, Top 30\% $\rightarrow$ 1x \\
GradientBoosting & ML & 30 & 60 & Top 30\% $\rightarrow$ 3x, Top 60\% $\rightarrow$ 1x \\
CNN\_LSTM & Hybrid & 30 & 60 & Top 30\% $\rightarrow$ 3x, Top 60\% $\rightarrow$ 1x \\
BiGRU & RNN & 15 & 50 & Top 15\% $\rightarrow$ 3x, Top 50\% $\rightarrow$ 1x \\
RandomForest & ML & 15 & 50 & Top 15\% $\rightarrow$ 3x, Top 50\% $\rightarrow$ 1x \\
NBEATS & DeepLearning & 10 & 20 & Top 10\% $\rightarrow$ 3x, Top 20\% $\rightarrow$ 1x \\
LightGBM & GradientBoosting & 10 & 50 & Top 10\% $\rightarrow$ 3x, Top 50\% $\rightarrow$ 1x \\
CatBoost & GradientBoosting & 5 & 30 & Top 5\% $\rightarrow$ 3x, Top 30\% $\rightarrow$ 1x \\
XGBoost & GradientBoosting & 15 & 20 & Top 15\% $\rightarrow$ 3x, Top 20\% $\rightarrow$ 1x \\
SeasonalNaive & TimeSeries & 5 & 40 & Top 5\% $\rightarrow$ 3x, Top 40\% $\rightarrow$ 1x \\
BiLSTM & RNN & 5 & 30 & Top 5\% $\rightarrow$ 3x, Top 30\% $\rightarrow$ 1x \\
DLinear & DeepLearning & 10 & 20 & Top 10\% $\rightarrow$ 3x, Top 20\% $\rightarrow$ 1x \\
NHiTS & DeepLearning & 15 & 20 & Top 15\% $\rightarrow$ 3x, Top 20\% $\rightarrow$ 1x \\
LSTM\_Attention & Attention & 15 & 50 & Top 15\% $\rightarrow$ 3x, Top 50\% $\rightarrow$ 1x \\
TFT & DeepLearning & 5 & 20 & Top 5\% $\rightarrow$ 3x, Top 20\% $\rightarrow$ 1x \\
AutoARIMA & TimeSeries & 5 & 20 & Top 5\% $\rightarrow$ 3x, Top 20\% $\rightarrow$ 1x \\
TCN & DeepLearning & 15 & 50 & Top 15\% $\rightarrow$ 3x, Top 50\% $\rightarrow$ 1x \\
Prophet & Prophet & 30 & 50 & Top 30\% $\rightarrow$ 3x, Top 50\% $\rightarrow$ 1x \\
ElasticNet & ML & 15 & 20 & Top 15\% $\rightarrow$ 3x, Top 20\% $\rightarrow$ 1x \\
Lasso & ML & 15 & 20 & Top 15\% $\rightarrow$ 3x, Top 20\% $\rightarrow$ 1x \\
ExponentialSmoothing & TimeSeries & 5 & 50 & Top 5\% $\rightarrow$ 3x, Top 50\% $\rightarrow$ 1x \\
GARCH & GARCH & 25 & 30 & Top 25\% $\rightarrow$ 3x, Top 30\% $\rightarrow$ 1x \\
Theta & TimeSeries & 5 & 20 & Top 5\% $\rightarrow$ 3x, Top 20\% $\rightarrow$ 1x \\
\bottomrule
\end{tabular}
\begin{tablenotes}
\small
\item Nota: $q_{3x}$ = percentil para posici\'on UPRO (3x). $q_{1x}$ = percentil para posici\'on SPY (1x). Predicciones por debajo del percentil $q_{1x}$ resultan en posici\'on Cash (0). Par\'ametros optimizados mediante grid search sobre conjunto de validaci\'on.
\end{tablenotes}
\end{table}


\subsection{An\'alisis de Par\'ametros}
\label{app:param_analysis}

\subsubsection{Patrones Observados}

\begin{enumerate}
    \item \textbf{Modelos conservadores} (AutoARIMA, TFT, BiLSTM): $q_{3x} = 5$, requieren alta confianza para apalancar.

    \item \textbf{Modelos agresivos} (GradientBoosting, Prophet, CNN-LSTM): $q_{3x} \geq 30$, m\'as dispuestos a apalancar.

    \item \textbf{Modelos equilibrados} (Ridge, BiGRU): $q_{3x} \in [15, 20]$, balance entre agresividad y conservadurismo.
\end{enumerate}

\subsubsection{Relaci\'on Par\'ametros-Rendimiento}

\begin{table}[htbp]
\centering
\caption{Correlaci\'on entre Par\'ametros y M\'etricas}
\label{tab:param_correlation}
\begin{tabular}{@{}lrr@{}}
\toprule
\textbf{Correlaci\'on} & \textbf{$q_{3x}$} & \textbf{$q_{1x}$} \\
\midrule
vs Retorno Total & 0.42 & 0.38 \\
vs Sharpe Ratio & 0.15 & 0.22 \\
vs Max Drawdown & 0.51 & 0.33 \\
vs \% Tiempo en 3x & 0.89 & 0.24 \\
\bottomrule
\end{tabular}

\vspace{0.2cm}
\small\textit{Nota: Correlaciones calculadas sobre los 23 modelos. $q_{3x}$ m\'as alto implica m\'as tiempo en UPRO, mayor retorno pero mayor drawdown.}
\end{table}

\subsubsection{Implicaciones}

\begin{itemize}
    \item \textbf{$q_{3x}$ alto + buen rendimiento}: El modelo tiene se\~nales fuertes frecuentes (GradientBoosting).
    \item \textbf{$q_{3x}$ bajo + buen rendimiento}: El modelo es selectivo pero preciso (Ridge).
    \item \textbf{$q_{3x}$ alto + mal rendimiento}: El modelo sobreestima se\~nales positivas (XGBoost).
\end{itemize}

\subsection{Sensibilidad a Par\'ametros}
\label{app:param_sensitivity}

\subsubsection{Ridge: An\'alisis de Sensibilidad}

\begin{table}[htbp]
\centering
\caption{Sensibilidad de Ridge a Par\'ametros de Umbral}
\label{tab:ridge_sensitivity}
\small
\begin{tabular}{@{}rrrrr@{}}
\toprule
\textbf{$q_{3x}$} & \textbf{$q_{1x}$} & \textbf{Retorno} & \textbf{Sharpe} & \textbf{Max DD} \\
\midrule
10 & 30 & +245.2\% & 0.87 & -32.1\% \\
15 & 30 & +287.6\% & 0.91 & -35.4\% \\
\textbf{20} & \textbf{30} & \textbf{+318.8\%} & \textbf{0.93} & \textbf{-38.6\%} \\
25 & 30 & +301.3\% & 0.85 & -42.3\% \\
20 & 40 & +298.4\% & 0.89 & -36.2\% \\
20 & 50 & +276.1\% & 0.84 & -34.8\% \\
\bottomrule
\end{tabular}

\vspace{0.2cm}
\small\textit{Nota: Valores en negrita indican par\'ametros \'optimos seleccionados.}
\end{table}

La sensibilidad moderada sugiere que los resultados son robustos a variaciones razonables en los par\'ametros.

\subsection{Estabilidad Temporal de Par\'ametros}
\label{app:param_stability}

\subsubsection{Walk-Forward Optimization}

Evaluamos si los par\'ametros \'optimos son estables en el tiempo usando walk-forward optimization:

\begin{enumerate}
    \item Dividir el per\'iodo de prueba en 5 sub-per\'iodos de ~260 d\'ias
    \item Re-optimizar par\'ametros al inicio de cada sub-per\'iodo
    \item Comparar par\'ametros seleccionados vs par\'ametros fijos
\end{enumerate}

\begin{table}[htbp]
\centering
\caption{Estabilidad de Par\'ametros: Ridge}
\label{tab:param_stability}
\begin{tabular}{@{}lrr@{}}
\toprule
\textbf{Sub-per\'iodo} & \textbf{$q_{3x}$ \'Optimo} & \textbf{$q_{1x}$ \'Optimo} \\
\midrule
Oct 2020 - Sep 2021 & 20 & 30 \\
Oct 2021 - Sep 2022 & 15 & 30 \\
Oct 2022 - Sep 2023 & 20 & 30 \\
Oct 2023 - Sep 2024 & 25 & 30 \\
Oct 2024 - Dic 2025 & 20 & 30 \\
\midrule
\textbf{Moda} & \textbf{20} & \textbf{30} \\
\bottomrule
\end{tabular}
\end{table}

Los par\'ametros \'optimos de Ridge son estables, con $q_{3x} = 20$, $q_{1x} = 30$ siendo \'optimos en 3 de 5 sub-per\'iodos.

\subsection{Interpretaci\'on de Umbrales}
\label{app:threshold_interpretation}

\subsubsection{Significado de los Percentiles}

Para Ridge con $(q_{3x}, q_{1x}) = (20, 30)$:

\begin{itemize}
    \item \textbf{Predicci\'on en top 20\%}: Posici\'on UPRO (3x leverage)
    \item \textbf{Predicci\'on en top 20-30\%}: Posici\'on SPY (1x)
    \item \textbf{Predicci\'on debajo de top 30\%}: Cash (tasa libre de riesgo)
\end{itemize}

Esto significa que Ridge solo toma posiciones de riesgo cuando su predicci\'on est\'a en el tercio superior de su distribuci\'on hist\'orica reciente.

\subsubsection{Percentiles vs Valores Absolutos}

El uso de percentiles rolling en lugar de umbrales fijos tiene ventajas:

\begin{enumerate}
    \item \textbf{Adaptabilidad}: Se ajusta autom\'aticamente a cambios en la distribuci\'on de predicciones.
    \item \textbf{Robustez}: Menos sensible a drift en el modelo.
    \item \textbf{Consistencia}: Mantiene proporci\'on constante de d\'ias en cada posici\'on (aproximadamente).
\end{enumerate}
